% Part: proof-theory
% Chapter: sequent-calculus
% Section: introduction

\documentclass[../../../include/open-logic-section]{subfiles}

\begin{document}

\olfileid{pt}{seq}{int}

\olsection{سريزه}

د سېکوېنټ حسابونه د !!{proof} نظامونه دي چې په 1930مو کلونو
کښې ګېرهرډ ګېنټسن (Gerhard Gentzen) معرفي کړل۔ موخه ئې تر
ډېره د شيانو د !!{proving} کولو عملي طريقه نۀ وه۔ بلکې ګېنټسن
دا د دې دپاره کارول چې د !!{proof}s د ممکنه جوړښت او خاصيتونو
په اړه ځينې پايلې ثابتې کړے شي۔ دا ډول پوښتنې کټ مټ هغه
پوښتنې دي چې د ثبوت تيوري ورسره سروکار لري۔

د سېکوېنټ حسابونه، لکه نوم ئې چې ښيي، پر \emph{سېکوېنټونو}
کار کوي، پر !!{formula}s نۀ۔ سېکوېنټ د !!{formula}s د دوو
لړيو يا څوګونو مجموعو يوه جوړه ده۔ د ګېنټسن اصلي نظامونو
(\Log{LK} او \Log{LJ}) لړۍ کارولې؛ اوس د ثبوت د تيورۍ
څېړونکي څوګونې مجموعې غوره ګڼي۔ څوګونې مجموعه د مجموعې
په شان ده، خو يو !!{element} له يو ځل څخه زيات لرلے شي۔
خو لکه په مجموعو کښې، د !!{element}s ترتيب اهميت نۀ لري۔
نو عبارتونه $\{!A, !A, !B\}$ او $\{!A, !B, !A\}$ هماغه يوه
څوګونې مجموعه بيانوي، خو دا څوګونې مجموعه له $\{!A,
!B\}$ او له $\{!A, !A, !B, !B\}$ څخه توپير لري، ځکه چې
په وروستنيو کښې د $!A$ او $!B$ شمېر د لومړۍ په نسبت بېل دے۔

په دوديزه توګه، د ثبوت په تيورۍ کښې د !!{formula}s څوګونې
مجموعې په لويو يوناني تورو ليکل کېږي۔ نو کله چې $\Gamma$، $\Delta$،
$\Pi$، $\Lambda$ يا $\Theta$ ليکو، موخه مو د هغه منطق د
!!{formula}s څوګونې مجموعې وي چې کار پرې کوو۔ همداراز، په
دوديزه توګه د ثبوت د تيورۍ څېړونکي د $\tuple{\Gamma, \Delta}$
د ليکلو پر ځاے دواړه برخې د سېکوېنټ په يوه نښه بېلوي؛
موږ~$\Sequent$ کاروو۔

\begin{defn}[سېکوېنټ]
\emph{سېکوېنټ} د دې بڼې يو عبارت دے:
\[
\Gamma \Sequent \Delta
\]
چې $\Gamma$ او $\Delta$ د !!{formula}s متناهي، او ښايي
تشې، څوګونې مجموعې دي۔ $\Gamma$ ته \emph{مقدم} او
$\Delta$ ته \emph{تالي} ويل کېږي۔
\end{defn}

فرض کړئ $!G$~په~$\Gamma$ کښې د !!{formula}s عطفي ترکيب دے
(تش عطفي ترکيب~$\ltrue$ ګڼو)، او $!D$~په~$\Delta$ کښې د
!!{formula}s فصلي ترکيب دے (تش فصلي ترکيب~$\lfalse$ ګڼو)۔
نو سېکوېنټ $\Gamma \Sequent \Delta$ په يو جوړښت~$\Struct{M}$
کښې رښتينے دے، که او يوازې که $!G \lif !D$ رښتينے وي۔
په کلاسيکي منطق کښې دا له دې خبرې سره يو شان ده چې يا
په~$\Gamma$ کښې د !!{formula}s څخه يوه دروغ وي، يا
په~$\Delta$ کښې د !!{formula}s څخه يوه رښتينې وي۔

که $\Gamma$ د !!{formula}s يوه څوګونې مجموعه وي، نو
$\Gamma, !A$ (يا $!A, \Gamma$) د هغې پايلې دپاره ليکو
چې $!A$ پر~$\Gamma$ ورزيات کړو، يعنې يوه څرګندېدنه ورزياتوو۔
که $\Delta$ هم د !!{formula}s يوه څوګونې مجموعه
وي، نو $\Gamma, \Delta$ د دواړو لړيو د څوګونو مجموعو اتحاد
دے؛ که $\Gamma$ د $!A$ د \emph{تکرار شمېر}~n سره ولري
(يعنې $n$ د~$!A$ نقلونه ولري)، او $\Delta$ ئې د تکرار
شمېر~$m$ سره ولري، نو $\Gamma, \Delta$ د $!A$ د تکرار
شمېر~$n+m$ سره لري۔ که د سېکوېنټ په يوه خوا کښې څوګونې
مجموعه تشه وي، عموماً هلته تش ځاے پرېږدو۔ د بېلګې په توګه،
$\Sequent !A$ هغه سېکوېنټ دے چې مقدم ئې تش دے او تالي
ئې $!A$ د تکرار شمېر~$1$ سره لري۔

!!^a{proof} د سېکوېنټ په حساب کښې د سېکوېنټونو يوه ونه ده:
تر ټولو پورته، يا ابتدايي، سېکوېنټونه په اړوند نظام کښې
بديهي اصول وي؛ او که يو سېکوېنټ له يوه يا دوو نورو سېکوېنټونو
لاندې وي، بايد په نظام کښې د استنتاج د يوې قاعدې له مخې
ترې په سم ډول راوتلے وي۔ لومړے د کلاسيکي منطق دپاره د
سېکوېنټ دوه بېلابېل نظامونه، $\Log{G1c}$ او~$\Log{G3c}$،
څېړو۔ بديهي اصول او قاعدې ئې په
\cref{pt:seq:int:tab:G1c,pt:seq:int:tab:G3c} کښې ورکړل شوي دي۔
\oliflabeldef{fol:seq::chap}{ (د ګېنټسن اصلي نظام \Log{LK}
په \olref[fol][seq][]{chap} کښې بيان شوے دے۔)}{}

\subfile{rules-G1c}
\subfile{rules-G3c}

نظامونه نازک توپيرونه لري، او له دې توپيرونو څخه راځي چې
د ثبوت د تيورۍ له مخې بېلابېل خاصيتونه لري۔ \Log{G1c} د
$!A
\Sequent !A$ بديهي اصول لري چې د نورو !!{formula}s اجازه
پکښې نۀ شته۔ \Log{G3c} د $!C, \Gamma \Sequent \Delta, !C$
په بڼه بديهي اصول لري چې په $\Gamma$ او~$\Delta$ کښې ئې
«څنګ» !!{formula}s هر ډول کېدے شي؛ خو په دواړو خواوو کښې
راغلے !!{formula}~$!C$ بايد \emph{اتومي} وي۔ \Log{G1c} د
$\LeftR{\land}$ او $\RightR{\lor}$ له قواعدو څخه د هرې يوې
دوه بڼې لري، خو \Log{G3c} د هرې يوې يوازې يوه بڼه لري۔
او \Log{G1c} اضافي «جوړښتي قاعدې» لري چې «انقباض»
(\LeftR{\Contraction}، \RightR{\Contraction}) او «کمزوري کول»
(\LeftR{\Weakening}، \RightR{\Weakening}) بلل کېږي۔

نظامونه \Log{G1c} او \Log{G3c} د کلاسيکي منطق دپاره سم
او بشپړ دي۔ په بله وينا، هر سېکوېنټ چې په دې نظامونو کښې
!!{provable} وي، په هر !!{structure} کښې رښتينے وي؛ او هر
سېکوېنټ چې په هر !!{structure} کښې رښتينے وي، !!a{proof}
لري۔ نو دا پوښتنه چې \emph{آيا} يو سېکوېنټ !!a{proof} لري،
د منطق په نورو طريقو، لکه د موډل د تيورۍ په طريقو، هم
ځوابېدې شي۔ او ځکه چې دواړه نظامونه کټ مټ هغه سېکوېنټونه
!!{prove} کوي چې په هر !!{structure} کښې رښتيني وي، نو په
ځانګړې توګه هماغه سېکوېنټونه ثابتوي۔

خو د ثبوت تيوري يوازې له دې سره سروکار نۀ لري چې \emph{آيا}
يو څۀ ثابتېدے شي، بلکې دا هم څېړي چې \emph{څنګه}۔ د ثبوت د
تيورۍ څېړونکي دا ډول پوښتنې څېړي: «آيا د يوې ځانګړې قاعدې
له کارولو څخه تل ډډه کېدے شي؟»، «آيا دا قاعده نظام ته داسې
ورزياتېدې شي چې نوي سېکوېنټونه د ثبوت وړ نۀ کړي؟»، يا
«آيا په يوه نظام کښې !!{proof}s په بل نظام کښې په !!{proof}s
بدلېدې شي؟» که د داسې يوې پوښتنې ځواب «هو» وي، دا حقيقت
اکثر د !!{proof}s پر پېچلتيا د استقرا په وسيله، يا په معادل
ډول د !!{proof}s پر جوړښت د استقرا په وسيله، ثابتېږي۔
له دې څخه يوازې د حقيقت ثبوت نۀ ترلاسه کېږي، لکه دا چې که
\Log{G1c} يو سېکوېنټ ثابتوي نو \Log{G3c} هم هماغه سېکوېنټ
ثابتوي، بلکې يوه کړنلاره هم ترلاسه کېږي، لکه په~\Log{G1c}
کښې !!{proof}s په~\Log{G3c} کښې په !!{proof}s بدلول۔

د ثبوت په تيورۍ کښې يوه مرکزي پايله \emph{د پرېکون لرې کولو
قضيه} ده۔ دا ښيي چې د پرېکون قاعده
\[
\Axiom$\Gamma \fCenter \Delta, !A$
\Axiom$!A, \Pi \fCenter \Lambda$
\RightLabel{\Cut}
\BinaryInf$\Gamma, \Pi \fCenter \Delta, \Lambda$
\DisplayProof
\]
د سېکوېنټ نظامونو ته ورزياتولے شو، پرته له دې چې دا بدل شي
چې کوم سېکوېنټونه !!{provable} دي۔ سربېره پر دې، ثبوت يوه
طريقه راکوي چې هر هغه !!{proof} چې د \Log{G1c} (يا \Log{G3c})
قواعد له \Cut{} سره يوځاے کاروي، په داسې ثبوت بدلوي چې
يوازې د \Log{G1c} (يا \Log{G3c}) قواعد کاروي۔ په بله وينا،
دا کړنلاره له هغو !!{proof}s څخه چې \Cut{} کاروي، دغه قاعده
«لرې کوي»؛ نوم ئې هم له همدې امله دے۔

\end{document}
